\scanpage{76}
Soit $T\subset X$ sous schéma, fini étale sur $S$, donc $T\simeq I_S$, $I$ ensemble fini, et soit $U=X\setminus T$ (donc $T$ est défini par une famille $(\unsure{q|g}_i)_{i\in I}$ de sections disjointes de $X$ sur $S$). \inserted{Si $g$ est le genre relatif, $\nu=\operatorname{card}I$, on supp. que $(g,\nu)\neq(0,1)$} Soit \deleted{$\unsure{f_0}$} $i_0\in I$, $f$ une section \deleted{disjointe}\inserted{de $X/S$ distincte des} $\unsure{q|g}_i$, et telle que \deleted{$\unsure{f_K(s)}=\unsure{q|g}_{i_0}$} $f$ et $\unsure{q|g}_{i_0}$ coïncident en $s$ (pt fermé de $S$). \deleted{Considérons} \inserted{Donc si $\eta=S\setminus\{s\}$} \deleted{et le \unsure{groupe de}} \deleted{par l'homomorphisme} on a donc un morphisme $\eta\longrightarrow U$, d'où $\pi_1(\eta)\longrightarrow\pi_1(U)$. Je dis que cet hom. n'est pas trivial, \deleted{\unsure{donc}} et $\unsure{\overline m|\overset{m}{=}}$, \inserted{si $\nu\geq2$,} que pour \inserted{\unsure{le dernier}}
\[
\diagram{\begin{tikzcd}[ampersand replacement=\&,column sep=small]
\& H_1(\eta,\mathbb Z_\ell) \arrow[r] \& H_1(U,\mathbb Z_\ell) \\
\inserted{T_\ell(\unsure{\overline K^*|\widehat K^*})} \arrow[ur,"\sim"] \& \& \inserted{[H_1(U_s,\mathbb Z_\ell)]} \arrow[u,"\unsure{\sim}"']
\end{tikzcd}}
\]
n'est pas trivial (pour, \deleted{tout} $\ell=$ nb premier distinct de la caractéristique résiduelle).

Comme la section rationnelle de $\mathcal J^1_{X/S}$ définie par $f$ est régulière, on voit que le composé de l'hom envisagé avec $H_1(\mathcal J^1_{X/S},\mathbb Z_\ell)$ est nul -- i.e. le $H_1(\eta,\mathbb Z_\ell)$ s'envoie dans la partie torique de $H_1(U,\mathbb Z_\ell)$ $\bigl(\xleftarrow{\sim}H_1(U_s,\mathbb Z_\ell)\bigr)$, qui est canoniquement isomorphe à $\deleted{\mathbb Z}\inserted{T}_\ell^I/\deleted{\mathbb Z}\inserted{T}_\ell$.

[NB cette partie est nulle si $\operatorname{card}I=1$, et
\margin{75}
